\par
\SourceLeaf{391}%
Es sei ferner \(r\). der Abstand des Mittelpunkts der gesuchten Ellipse von dem Anfangspunkte der Abscissen, und \(\varphi\) die Neigung der von letzterm zu ersterm gezogenen geraden Linie gegen die Abscissen-Linie, oder \(r\cos\varphi,r\sin\varphi\) die Coordinaten des Mittelpunkts der Ellipse. Man findet hieraus leicht, dass das Perpendikel von diesem Mittelpunkte auf die erste Seite des Vierecks
\[
=a-r\cos(A-\varphi)
\]
sein werde; auf ähnliche Art werden die Perpendikel auf die drei andern Seiten ausgedrückt.

Bezeichnet man die halbe grosse Axe der Ellipse mit \(\alpha\), die halbe kleine Axe mit \(\beta\), die Neigung der letztern gegen die Abscissen-Linie mit \(\psi\), so ist offenbar \(A-\psi\) die Neigung des Perpendikels aus dem Mittelpunkte auf die erste Seite des Vierecks gegen die kleine Axe, welches, wenn jene die Ellipse berühren soll, nach bekannten Gründen durch
\[
\BareRoot{[\alpha\alpha\sin(A-\psi)^2+\beta\beta\cos(A-\psi)^2]}
\]
ausgedrückt wird. Man hat also die Gleichung
\[
a-r\cos(A-\varphi)=\BareRoot{[\alpha\alpha\sin(A-\psi)^2+\beta\beta\cos(A-\psi)^2]}
\]
und eben so drei andere ganz ähnliche, wenn man statt \(a\) und \(A\) die sich auf die andern Seiten beziehenden Zeichen substituirt. Schafft man also die Irrationalität weg, und setzt Kürze halber
\[
\begin{aligned}
rr-\alpha\alpha-\beta\beta&=t\\
\alpha\alpha-\beta\beta&=u
\end{aligned}
\]
so sind unsere vier Gleichungen
\begingroup\fontsize{10}{12}\selectfont
\[
\begin{array}{rl}
\mathrm{I.}&2aa+t-4ar\cos(A-\varphi)+rr\cos2(A-\varphi)-u\cos2(A-\psi)=0\\
\mathrm{II.}&2a'a'+t-4a'r\cos(A'-\varphi)+rr\cos2(A'-\varphi)-u\cos2(A'-\psi)=0\\
\mathrm{III.}&2a''a''+t-4a''r\cos(A''-\varphi)+rr\cos2(A''-\varphi)-u\cos2(A''-\psi)=0\\
\mathrm{IV.}&2a'''a'''+t-4a'''r\cos(A'''-\varphi)+rr\cos2(A'''-\varphi)-u\cos2(A'''-\psi)=0
\end{array}
\]
\endgroup
Multiplicirt man die erste Gleichung mit \(\sin2(A''-A')\), die zweite mit \(\sin2(A-A'')\), die dritte mit \(\sin2(A'-A)\), und addirt die Producte, so wird (m.\,s.\ Art.\,78 meiner \textit{Theoria motus corporum coelestium})
